\documentclass[11pt]{article}
\usepackage[margin=1in]{geometry}
\usepackage{amsmath,amssymb,amsthm,graphicx,microtype}
\usepackage[colorlinks=true,linkcolor=blue,urlcolor=blue,citecolor=blue]{hyperref}
\newtheorem{theorem}{Theorem}
\newtheorem{proposition}{Proposition}
\newcommand{\R}{\mathbb R}
\newcommand{\Vone}{V_1}
\title{A two-dimensional spiral-dilation counterexample\\
\large to Conjecture 2.14 of arXiv:1812.07375}
\author{Autonomous proof audit --- lane 9}
\date{17 August 2026}

\begin{document}
\maketitle

\begin{center}
\fbox{\parbox{0.91\linewidth}{\textbf{Disposition.}
The conjecture is false, but the negative answer is already present in the
literature: Moisala's 2020 thesis gives a geodesic $L^1$ example, and
Moisala--Pasqualetto explicitly point to it.  The proof below records a simpler
finite-dimensional counterexample obtained from the rotation template of
Moisala--Pasqualetto.  No novelty claim is made.}}
\end{center}

\section{Source question and transcription}

The source paper asks (Conjecture 2.14, printed page 12):
\begin{quote}
If $G$ is a metric scalable group that is separable and nilpotent and its
distance is geodesic, then its first layer $\Vone(G)$ generates $G$.
\end{quote}
Here a scalable group has automorphisms $\delta_t$ for $t\ne0$, with
$\delta_t\delta_s=\delta_{ts}$, a continuous map $(t,x)\mapsto\delta_t x$,
and $\delta_0x=e$.  A compatible left-invariant metric satisfies
\[
 d(\delta_t x,\delta_t y)=|t|d(x,y).
\]
The first layer is
\[
 \Vone(G)=\{x:\delta_{t+s}x=\delta_t x\,\delta_s x
                  \text{ for every }s,t\in\R\}.
\]
``Generates'' means that the closed subgroup generated by its dilates is all
of $G$.

\begin{figure}[ht]
\centering
\includegraphics[width=0.94\linewidth]{figures/open_problem_crop.png}
\caption{Source-page evidence: arXiv:1812.07375, printed page 12.}
\end{figure}

\section{Idea and proof intuition}

Metric homogeneity determines the radial size of a dilation but not its
angular part.  Let the angular part rotate logarithmically with the scale.
The logarithm converts multiplication of parameters into addition of angles,
so the scalable-group law remains true.  The orbit of any nonzero vector then
spirals as it contracts.  Such an orbit cannot be a one-parameter subgroup
for vector addition, and in fact the single test $s=t=1$ already excludes
every nonzero vector from the first layer.

\section{Counterexample}

For $a\in\R$, let $R_a$ be counterclockwise rotation of $\R^2$ through angle
$a$.  Take $G=(\R^2,+)$, the Euclidean metric, and define
\[
 \delta_0x=0,\qquad
 \delta_t x=tR_{\log|t|}x\quad(t\ne0).
\]

\begin{theorem}
The triple $(G,\delta,d_2)$ is a complete separable geodesic metric scalable
group.  Its underlying group is abelian, hence nilpotent, but
$\Vone(G)=\{0\}$.  Consequently $\Vone(G)$ does not generate $G$.
\end{theorem}

\begin{proof}
For $t\ne0$, $\delta_t$ is an invertible linear map and therefore a group
automorphism.  If $s,t\ne0$, commutativity of planar rotations and
$\log|st|=\log|s|+\log|t|$ give
\begin{align*}
 \delta_t(\delta_sx)
 &=tR_{\log|t|}\bigl(sR_{\log|s|}x\bigr)\\
 &=stR_{\log|st|}x=\delta_{st}x.
\end{align*}
The identities with $st=0$ follow from $\delta_0x=0$.  For $t\ne0$,
\[
 \|\delta_t x-\delta_t y\|_2=|t|\|x-y\|_2,
\]
because rotations are isometries.  This also proves continuity at $t=0$:
if $(t_n,x_n)\to(0,x)$, then
$\|\delta_{t_n}x_n\|_2=|t_n|\|x_n\|_2\to0$.
Joint continuity away from zero is immediate from the formula.  Thus all
scalable-group and metric-compatibility axioms hold, including for negative
parameters.

Euclidean $\R^2$ is complete, separable, and geodesic, and its additive group
is abelian.  Finally, if $x\in\Vone(G)$, the defining identity with $s=t=1$
would imply
\[
 2R_{\log2}x=\delta_2x=\delta_1x+\delta_1x=2x.
\]
Since $0<\log2<2\pi$, the nontrivial rotation $R_{\log2}$ has no nonzero
fixed vector.  Hence $x=0$.  Conversely $0\in\Vone(G)$, so
$\Vone(G)=\{0\}$ and its generated closed subgroup is proper.
\end{proof}

\section{Literature status and verification}

The original conjecture is already answered negatively.  On page 21 of the
introductory part of her thesis, Moisala constructs a geodesic example on the
additive group $L^1(\R)$ with
$(\delta_\lambda f)(x)=f(x/\lambda)$ and proves that its first layer is
trivial.  Moisala--Pasqualetto later state explicitly that an abelian complete
metric scalable group with trivial first layer ``can even be made geodesic'',
citing that thesis page.

Their Remark 3.4 also uses the general rotation template
$\delta_t=(tI)R_{f(t)}$ for a multiplicative-to-additive homomorphism $f$.
Choosing the continuous homomorphism $f(t)=\log|t|$ gives precisely the
counterexample above.  The argument has been checked against every axiom in
the source definition; the test $s=t=1$ independently certifies the trivial
first layer.  Local compactness causes no contradiction: metric
characterizations may identify the underlying metric group as a Carnot group
without forcing an arbitrarily specified scalable structure to be its
standard Carnot dilation family.

\paragraph{Limitations and novelty.}
This packet refutes only the printed conjecture.  It does not classify which
extra compatibility condition on a given dilation family forces generation
by the first layer.  The negative answer is prior literature.  The
two-dimensional presentation is an independently derived simplification of a
published rotation template, not asserted to be new.

\begin{thebibliography}{9}
\bibitem{LDLM}
E. Le Donne, S. Li, and T. Moisala,
\emph{G\^ateaux differentiability on infinite-dimensional Carnot groups},
arXiv:1812.07375, Conjecture 2.14 (2018).

\bibitem{Moisala}
T. Moisala,
\emph{Unraveling Intrinsic Geometry of Sets and Functions in Carnot groups},
JYU Dissertations 264, University of Jyv\"askyl\"a (2020), introductory p.~21,
\url{https://jyx.jyu.fi/handle/123456789/71351}.

\bibitem{MP}
T. Moisala and E. Pasqualetto,
\emph{Direct limits of infinite-dimensional Carnot groups},
arXiv:2101.03979, Remarks following Proposition 4.5 and Remark 3.4 (2021).
\end{thebibliography}

\end{document}
